\section{Discussion}

\subsection{Calibration and cross-validation: Inference of PMoAs from life-history data}
\noindent
By fitting alternative PMoAs to effect data for larvae, 
we could draw some conclusions regarding the plausibility of PMoAs. 
For a given effect on growth, 
PMoAs $M$ and $A$ are associated with a clear delay of metamorphosis, 
whereas PMoA $G$ only implies a subtle delay of metamorphosis. 
This can be explained by the interaction of these PMoAs with mass fluxes in the DEB model.\\
PMoA $G$, i.e. a decrease in growth efficiency, 
only causes an indirect effect on the maturation rate.
Furthermore, any reduction in growth also reduces maintenance costs, 
and the effect of $G$ approaches zero as an individual reaches its maximum length 
(the growth efficiency $\eta_{AS}$ does not affect the equilibrium conditions of the model).
PMoAs $M$ and $A$ on the other hand have a direct effect on growth and maturation.\\
PMoA $\kappa$ represent a special case, as it does not imply a loss of energy, 
but only a change in the allocation of energy. 
In the present study, $\kappa$ did not explain observed effects. 
Generally, a decrease in $\kappa$ might be a plausible candidate if premature metamorphosis is observed 
in combination with adverse effects on growth.\\
The prediction of effects at GS 46 depended heavily on the choice of PMoA. 
The message here is two-fold: 
At the one hand, extrapolations from larval effect data to GS 46 should be done with caution. 
We suggest a model averaging approach if the PMoA identification from larval data is ambiguous.
At the other hand, effect data at GS 46 apparently carries a lot of information regarding the PMoA. 
Generating such effect data is resource-intensive, but for cases where an unequivocal PMoA identification is crucial, 
it could be worth the effort. 
%
\subsection{Validation: Predicting mixture effects}
\noindent
The dominance effect of 2,4-D over Flupyradifurone could be predicted with high accuracy. 
Given that the substances acted via differnt PMoAs, a diferentiation between IA and DA models was obsolete. 
The presence of a dominance effect can be explained through the combination of PMoAs, 
based on the same rationale that was used to select the PMoAs in the first place. 
We identified $G$ as most likely PMoA for Flupyradifurone, 
implying that effects weaken over time and only indirect effects on maturation. 
We identified $M$ as most likely PMoA for 2,4-D, 
implying effects across the entire growth curve, as well as direct effects on the maturation rate. 
%
\subsection{Implications for ampibian population dynamics}
\noindent
While the present study focussed on the modelling of sublethal effects on the individual level, 
one of the goals of the amphibian DEB-TKTD model development is to extrapolate effects to the population level.
Achieving this would at least require that the present analysis is extended to also account for lethal effects. 
In fact, the severity of lethal effects is highly life stage-specific (Gonzalez-Lopez et al., unpublished), 
with the highest mortality rate ocurring during metamorphosis, 
manifesting metamorphosis as a bottleneck in the amphibian life history.\\ 
Apart from direct lethal effects, the ecology of amphibians implies that sublethal effects in laboratory tests 
may translate to letahl effects in the field.\\
In ephemeral wetlands such as those inhabited by \textit{D. galganoi},
a delay in metamorphosis could lead to death if individuals fail to reach metamorphosis 
before the aquatic habitat has dried out. 
This means that the results presented in the present study could lead 
to dramatically different effects on population dynamics if the effect of drying is considered. 
Simultaneously, it is known that amphibians may accelerate development under drying conditions, 
counter-acting the effects described here. 
Such acceleration of development may in turn lead to adverse effects in later life stages. 
These are relationships between chemical exposure and amphibian life-history that have not been considered 
in the present study, but could be integrated into a DEB-TKTD modelling approach.
%
\subsection{Conclusions}
\noindent 
We have demonstrated that DEB-TKTD models can be used to perform predictive modelling 
of pesticide effects in amphibian larvae and metamorphs.
The extrapolation of effects across life stages depends heavily on the selection of the PMoA. 
DEB-TKTD models can be fitted to larval data up to GS 42 and some inference about the PMoA can be made, 
but data on effects at GS 46 is highly informative regarding the PMoA.
Generally, the model has a tendency to under-predict effects on the mass at metamorphosis. 
The effects of a binary mixture of pesticides with different PMoAs could be predicted with high accuracy.